\label{sec:DescartesNightmare}
\begin{quotex}
We live among ruins in a World in which ``god is dead'' as Nietzsche stated. The ideals of today are comfort, expediency, surface knowledge, disregard for one's ancestral heritage and traditions, catering to the lowest standards of taste and intelligence, apotheosis of the pathetic, hording of material objects and possessions, disrespect for all that is inherently higher and better - in other words a complete inversion of true values and ideals, the raising of the victory flag of ignorance and the banner of degeneracy. In such a time, social decadence is so widespread that it appears as a natural component of all political institutions. The crises that dominate the daily lives of our societies are part of a secret occult war to remove the support of spiritual and traditional values in order to turn man into a passive instrument of dark powers. The common ground of both Capitalism and Socialism is a materialistic view of life and being.

\flright{\textsc{S. H. Nasr}}

The rule and criterion of truth is to have made it. Hence the clear and distinct idea of the mind not only cannot be the criterion of other truths, but it cannot be the criterion of that of the mind itself; for while the mind apprehends itself, it does not make itself, and because it does not make itself it is ignorant of the form or mode by which it apprehends itself

\flright{\textsc{Giambattista Vico} (commenting on Descartes)}

All that is solid melts into air, all that is holy is profaned, and man is at last compelled to face with sober senses, his real conditions of life, and his relations with his kind.

\flright{\textsc{Marx and Engels}, \textit{The Communist Manifesto}}

\textit{La lotta differenzia, seleziona, crea gerarchia; soprattutto quando – per usare delle espressioni tradizionali – non è la piccola lotta, ma la grande lotta; non la lotta di uomo contro uomo or contro l'ambiente, ma la lotta dell'elemento sovrannaturale dell'uomo contro tutto ciò che in lui è natura, sensazione, materialità, agitatizione, miraggio di vana grandezza; contro il caos che è in lui, primi di esser fuor di lui.}

The battle differentiates, selects, creates hierarchy; especially when – to use the traditional terms – it is not the lesser battle, but the greater battle; not the battle of man against man, or against the world, but the battle of the supernatural element of man against everything in him that is nature, sensation, materiality, ferment, delusion of grandeur; against the chaos that is within him, before it is outside him.

\flright{\textsc{Julius Evola}}

Taluni tornano sui loro passi, che nuove forze sorgano per la riconquista.

\flright{\textsc{Julius Evola}}
\end{quotex}

\paragraph{Introduction} The fundamental difference between the traditional and modern world views is the knowledge of a spiritual level of reality that transcends our ordinary experience. This spiritual level is the source of all manifestation and is known, not through empirical evidence or a process of reasoning, but rather through a direct “knowing”, called Intuition or gnosis. This gnosis is not irrational, but is supra-rational, that is, above reason, though not contrary to reason.

The traditional teaching is that man is a tripartite being, that is, a Unity of Spirit, soul, and body. These states of being are hierarchically arranged, with the lower state subordinated to the higher states.

The modernist denies the existence of a spiritual state of being, and thus, of any knowledge that transcends the senses or the reasoning faculty. Science ignores it completely by its very methodology, and regards any claim to the contrary as delusional. The major exoteric religions accept the existence of the spiritual realm, but insist it can be known only by faith, with a caste of priests or ministers keeping guard over it.

Now, the difference between the Traditional and modern world view must not be understood in a temporal sense, but rather ontologically. “Tradition” indicates a transcendental source that impinges on the empirical at every time and place. “Modern” indicates the denial of such transcendence, viewing the empirical as self-contained. Thus, the Traditional and modern world views are independent of time and place, and have always been with us, with one predominating over the other. Other times, these two spirits are found battling each other for influence.

In \textit{Imperialismo Pagano}, Evola applies these ideas to the particular historico-political situation he found himself in. Opposed to any sort of revolutionary impulse, Evola hoped to draw Italy back to its spiritual roots in the ancient Roman Imperium. He felt that that spirituality was more “natural” and congenial to the people of Italy. Evola analyzes in exquisite detail the root causes and results of the spiritual warfare between the discordant world views then tearing Europe apart. Evola traces certain viewpoints to their sources in various religions, nations, and ethnic groups. It is important to keep in mind that the battle is first and foremost within one's own consciousness; the outer world is a symptom or result of particular mind sets, or their admixture.

Evola ranges over topics including mythology, philosophy, history, and political science. His ideas challenge us and force us to see beyond our everyday patterns of thought that we often absorb without reflection. In the end, no matter which side of the issue you end up on, you will know precisely why you lean toward one side rather than the other.

\paragraph{Modernism} Descartes did not originate modernism, but he outlined its essential features; thus, he serves as a useful foil to delineate its essential features. In contrast with the unity of the tripartite being, Descartes introduces “dualism” – that is, the concept that thought and matter are distinct substances, with no direct relation between the two. Denying any level (spirit) higher than thought and denying any form of knowing (gnosis, intuition) higher than reason, Descartes elevates thought to the pinnacle of human life. It becomes man's very identity – “I think, therefore I am”. A man is identified with the stream of thought he experiences. This leads to some anti-traditional consequences. If I am my thoughts, and there is no higher knowledge, then mere opinion becomes the standard of truth. To deny my opinions is to deny my self; and to deny a higher knowledge means that my opinions are as good as the next fellow's.

For the traditionalist, the “I am” (Spirit) transcends both thought (Soul) and matter (Body). As the Observer, the Spirit is detached both from the stream of thought and the concupiscence of the body. As the Unmoved Mover, the Spirit the source of all creative activity without itself being part of that activity.

This leads to stability and continuity. Detached from the vagaries of public opinion, the traditionalist holds firm to eternal truths. The modern mind can be swayed by the “winds of doctrine”, by the opinion “du jour”. Thus we see wild changes among the masses, based solely on transient events or mere collective whims.

The second key aspect of the Cartesian system is the method of doubt. Descartes proposed to doubt everything until only “clear and distinct” ideas remained. Since the Traditional teachings transcend reason, they cannot be understood by the modern mind, who can only see in them error, deception, the “opiate of the masses”. Everything arising out of the life of the Spirit is subject to radical doubt – family, nation, race, religion, gender, marriage, etc. Knowledge is replaced by opinion. Public opinion is enforced through rhetoric, “gerede”, the idle chatter of the mass man, who becomes a consumer of information. Where that fails, recourse is made to force and violence to enforce thought within narrow confines. In contrast, the traditional view relies on persuasion – the measure of one's strength is not force, but rather the impact of “power-ideas” to move men's minds.

The third element of dualism is that the body is separated from the mind, and so it is no longer subordinated to anything higher. Therefore, the needs and desires of the body are of fundamental – and only – concern. There is no higher spiritual force directing the body – therefore everything is reduced to the material, biological, evolutionary, economic scale.

The whole nature of personhood is altered. No longer are there persons – in the sense of a center of consciousness and will – that are hierarchically arranged and organically related. Instead, there are “individuals”, independent of each other, all presumably equal and identical. Their relations are not organic, but rather contractual – they are entered into as a matter of convenience to the parties involved and have no higher purpose or meaning. An individual is not his own center, or end in himself ... he is a consumer of outside forces: of ideas, of entertainment, of goods and services, all in service to the needs and desires of the body. Anything that thwarts the satisfaction of desires is “bad” – not just bad, but absurd, beyond understanding. The Person of tradition who aspires to an heroic and ascetical life is considered a fool, or a “troglodyte”, on the way to extinction.

Since the needs and desires of the body are the least common denominator, every individual is equal to every other. As the quest for “clear and distinct” ideas if futile, all that remain are random thoughts, so one man's opinion is as good as another. Since there is no higher standard, the only people worth emulating or respecting are movie star, the businessman, the sportsman.
